\section{Current Situation}

\subsection{Problem Statement}

The multilateral trading system was built on the assumption that economic interdependence would effectively contribute to greater international stability and prosperity, as countries would face higher downsides to waging wars and causing disruptions in global supply chains. Since the establishment of the General Agreement on Tariffs and Trade in 1947 and the subsequent creation of the World Trade Organization in 1995, states have largely accepted the rules-based framework designed to reduce barriers to trade and prevent discriminatory economic practices. Nonetheless, the contemporary geopolitical environment has increasingly challenged this model, as growing strategic competition between major powers, the resurgence of industrial policy and heightened concerns regarding economic security have blurred the distinction between trade policy and national security.

At the center of this debate lies a fundamental dilemma positing that while governments have a legitimate interest and responsibility to protect their national security, the increasing use of security-based trade measures risks undermining the openness, predictability and non-discrimination upon which the WTO is founded. This fragmentation of global trade is occurring while growing concerns about the distribution of benefits of globalization exist. While the expansion of global trade contributed to unprecedented economic growth, critics have long argued that its benefits have been unevenly distributed both within and between nations. Joseph Stiglitz in Globalization and Its Discontents, argued that the governance of globalization through organizations such as the World Bank, the International Monetary Fund, or the World Trade Organization, often failed to account for the interests of developing economies and vulnerable populations. Consequently, concerns regarding the fragmentation of trade intersect not only with matters of security but also with longstanding interrogations on equity and inclusion within the global economic system. The key question to confront is therefore to understand how states can address genuine security and equity concerns without eroding the foundations of free trade.

\subsubsection{The Expansion of Security-Based Trade Measures}

The growing use of national security as a justification for trade restrictions has become a defining trend of contemporary international economic relations, as the GATT’s Article XXI permits WTO members to adopt measures they consider necessary for the protection of their essential security interests. Nevertheless, this provision was rarely invoked and was generally viewed as an exceptional safeguard reserved for extraordinary circumstances.

This provision became less exceptional with governments increasingly relying on national security arguments to justify their tariffs, exports controls, investments restrictions and industrial subsidies.\footnote{Van Den Bossche, P., Akpofure, S., \& World Trade Institute. (2020). The Use and Abuse of the National Security Exception under Article XXI(b)(iii) of the GATT 1994 The Use and Abuse of the National Security Exception under Article XXI(b)(iii) of the GATT 1994. \url{https://www.wti.org/media/filer_public/50/57/5057fb22-f949-4920-8bd1-e8ad352d22b2/wti_} working\_paper\_03\_2020.pdf} A notable case illustrating this phenomenon occurred in 2018 when the United States of America (USA) imposed tariffs on imported steel and aluminium under Section 232 of the Trade Expansion Act (US Congress, 2025) . The primary argument to back such a decision was that dependence on foreign suppliers could threaten the national security of the USA. Several WTO member states challenged the measures by arguing that they constituted disguised protectionist policies rather than legitimate security actions.\footnote{European Commission. (2018, May 31). Press corner. European Commission - European Commission. \url{https://ec.europa.eu/commission/presscorner/detail/en/IP_18_4006}} Therefore, a question arises, if states possess unlimited discretionary powers giving them the ability to define what constitutes a security threat, WTO obligations may become very difficult to enforce.

Governments are now treating supply chains, energy resources, critical minerals, telecommunications infrastructure, and advanced technologies as strategic assets (OECD, 2024). As a result, economic policy has become increasingly intertwined with national security policy, creating new tensions between sovereign decision-making and multilateral trade commitments.

\subsubsection{Paralysis of the WTO Dispute Settlement System}

The challenge that poses expanding security-based measures has been multiplied by the growing difficulties the WTO’s dispute settlement system faces. Indeed, since December 2019, the WTO Appellate Body has remained non-functional due to disagreements among member states regarding its role and ultimately its authority (European Commission, 2023). As geopolitical tensions increased, governments have become more willing to act unilaterally rather than relying on a paralysed body, doubting the capacities of WTO mechanisms to effectively address their issues.

\subsubsection{Fragmentation of Global Trade Networks}

The increasing securitization of trade has contributed to a broader restructuring of global economic relationships, as for decades globalization was primarily guided by considerations of efficiency, cost reduction and comparative advantage. Nowadays, other elements are added to the equation, with governments increasingly prioritizing resilience and strategic autonomy.

In this context, concepts such as friendshoring and nearshoring have become crucial to economic policymaking and global trade. Friendshoring refers to the relocation of production and supply chains toward politically aligned partners, while nearshoring emphasizes the geographical proximity as a means of reducing vulnerability and costs.\footnote{Nedumpara, J. (2024, March 31). Friendshoring, Nearshoring, Greenshoring and Resshoring: Changing Faces of Global Supply Chains and its Impact on International Economic Law Introduction to the Special Issue. WTO Chairs Programme. \url{https://wtochairs.wto.org/friendshoring-nearshoring-greenshoring-and-resshoring-changi} ng-faces-global-supply-chains-and-its} While such strategies may improve resilience against disruptions, they also risk fragmenting the global economy into competing economic blocs.\footnote{Crowe, D., \& Rawdanowicz, Ł. (2023). Risks and opportunities of reshaping global value chains. OECD Economics Department Working Papers. \url{https://doi.org/10.1787/f758afe8-en}} The semiconductor sector provides one of the clearest examples of this phenomenon. Advanced semiconductors play a critical role in artificial intelligence, telecommunications, defense systems, and advanced manufacturing. Since 2022, the United States has progressively expanded export controls limiting China's access to advanced semiconductor technologies, citing national security concerns. In response, China has implemented its own restrictions on the export of critical minerals and accelerated efforts to achieve technological self-sufficiency.\footnote{Branstetter, L. (2024, July). EXPORT CONTROLS AND US-CHINA TECHNOLOGY COMPETITION IN AN INTERDEPENDENT WORLD. Economic Studies at Brookings. \url{https://www.brookings.edu/wp-content/uploads/2024/07/20240701_Branstetter_Sanctions}. pdf} These developments illustrate a broader shift in international trade. Increasingly, states are willing to accept economic costs in pursuit of strategic objectives. While such measures may enhance national resilience, they also challenge the WTO's traditional commitment to open and non-discriminatory trade.

\subsubsection{Unequal Consequences for Developing Economies}

While much of the debate has been focused on major powers, the consequences of trade fragmentation are often most significant for developing countries relying heavily on access to global markets, foreign investment and their active role in international supply chains. As advanced economies are redirecting investment and production toward politically aligned partners instead of maximizing profits and considering comparative advantages, some developing countries may struggle to attract capital or integrate into strategic industries. On one hand, some may face pressure to align with one geopolitical bloc or another, potentially limiting their economic and diplomatic flexibility. On the other hand and at the same time, many developing countries can not mobilize the necessary financial resources to implement large-scale industrial policies comparable to those adopted by major powers. Nevertheless, fragmentation may also create opportunities. Some countries have benefited from the relocation of manufacturing and supply chains away from traditional production hubs. States such as Vietnam, Indonesia, and Mexico have attracted increased investment as firms seek to diversify production networks. This demonstrates that the effects of geopolitical competition are neither uniform nor entirely negative.

\subsection{Possible Solutions}

The debate surrounding national security and free trade ultimately raises broader questions regarding the future of multilateralism. The extensive sanctions imposed following Russia's invasion of Ukraine, the growing strategic rivalry between the United States and China, and the European Union's pursuit of a “de-risking” strategy all demonstrate that governments increasingly view trade through a geopolitical lens.

Nonetheless, these developments have not eliminated the need for international cooperation as global supply chains remain deeply interconnected and many contemporary challenges such as climate change, technological governance or international economic development require coordinated and global efforts. The challenge facing the WTO is therefore not whether security concerns should influence trade policy, but rather how such concerns can be managed without undermining the rules and institutions that facilitate international commerce. Delegates should grapple with two interconnected questions: how can security concerns be addressed without destroying free trade, and how can the WTO remain relevant in an era increasingly defined by great-power competition?

Given the complexity of the issue at hand, it is not expected that one single solution will resolve all tensions between the two sides of the debate. Nevertheless, several policy pathways have emerged within academic, governmental and multilateral discussions, each attempting to preserve a balance between state sovereignty and the integrity of the WTO system.\footnote{Maruyama, W., \& Wolff, A. Wm. (2023). Saving the WTO from the national security exception. Peterson Institute for International Economics.} \footnote{WTI Working Paper No. 03/2020, 2020}

One approach focuses on clarifying the interpretation of GATT Article XXI as many scholars argued that the provision is not meant to be a completely unconstrained “self-judging” exception and tWTI Working Paper No. 03/2020, 2020hat its use should remain subject to good-faith and some level of objective criteria.\footnote{Ikeda, T. (2023). A proposed interpretation of GATT Article XXI(b)(ii) in light of its jurisprudence and practice. Cornell International Law Journal.} In this view, clearer guidance on the scope of “essential security interests,” the meaning of “necessity,” and the types of circumstances covered by the exception could reduce legal uncertainty while preserving states’ ability to respond to genuine threats.\footnote{Meti. (2020). COLUMN: SECURITY EXCEPTIONS -INTERPRETATION OF GATT ARTICLE XXI. \url{https://www.meti.go.jp/english/report/data/2020WTO/pdf/column_02.pdf}} Such clarification would not eliminate the national security exception, but it could make it harder to invoke Article XXI as a blanket justification for trade protectionism.\footnote{Mavroidis, P. C. (n.d.). Litigating national security in the WTO era. Columbia Law School Scholarship Archive. \url{https://scholarship.law.columbia.edu/faculty_scholarship/4656/}}

A second proposal involves reforming and strengthening the WTO dispute settlement system. The institutional weakness created by the Appellate Body crisis has made it increasingly difficult for members to rely on it, especially on politically sensitive disputes.\footnote{European Parliament, \& Grieger, G. (2024, June). World Trade Organisation Appellate Body crisis and the multi-party interim appeal arbitration arrangement. \url{https://www.europarl.europa.eu/RegData/etudes/BRIE/2024/762342/EPRS_BRI(2024)76234} 2\_EN.pdf} Restoring a fully functioning appellate mechanism or designing specialized procedures for security-related disputes could improve confidence in the WTO’s enforcement capacity and reduce incentives for unilateral retaliation.\footnote{CSIS. (2025). The World Trade Organization: The Appellate Body Crisis | Economics Program and Scholl Chair in International Business | CSIS. Csis.org. \url{https://www.csis.org/programs/economics-program-and-scholl-chair-international-busine} ss/world-trade-organization} \footnote{Kommerskollegium. (2023). The WTO Appellate Body Crisis. \url{https://www.kommerskollegium.se/globalassets/publikationer/rapporter/2023/the-wto-app} ellate-body-crisis.pdf} As several scholars note, the problem is not only that national security measures are difficult to assess, but that the current system lacks the institutional credibility needed to resolve them consistently.\footnote{Maruyama, W., \& Wolff, A. Wm. (2023). Saving the WTO from the national security exception. Peterson Institute for International Economics.}

Other proposals emphasize transparency as a key aspect to reconcile the two elements, indeed, they argue that enhanced notification requirements, consultation mechanisms and institutionalized information-sharing practices could improve trust among states and reduce the risk that security measures escalate into broader trade conflicts. Greater transparency may also help distinguish legitimate security measures from protectionist policies by forcing states to explain the rationale, scope, and proportionality of their actions more clearly.\footnote{Ikeda, T. (2023). A proposed interpretation of GATT Article XXI(b)(ii) in light of its jurisprudence and practice. Cornell International Law Journal.} In this sense, transparency is not a substitute for legal reform, but a practical complement to it: even if members disagree on how far Article XXI should extend, they may still support stronger procedural obligations that make abuse easier to detect and harder to justify.\footnote{Mavroidis, P. C. (n.d.). Litigating national security in the WTO era. Columbia Law School Scholarship Archive. \url{https://scholarship.law.columbia.edu/faculty_scholarship/4656/}}
